\documentclass[12pt, ukrainian]{article}
\usepackage[a4paper]{geometry}
\setlength\parindent{0mm}
\setlength\parskip{4mm}
\geometry{top=1.5cm,bottom=1.5cm,left=1.3cm,right=1.5cm}
\pagestyle{empty}
\usepackage[T2A]{fontenc}
\usepackage[utf8]{inputenc}
\usepackage[ukrainian]{babel}
\usepackage{graphicx}
\usepackage{caption}
\usepackage{caption}
\DeclareCaptionFormat{nocaption}{}
\captionsetup{format=nocaption,aboveskip=0pt,belowskip=0pt}
\usepackage[Export]{adjustbox} 
\adjustboxset{max size={0.9\linewidth}{0.9\paperheight}}
\usepackage{verbatim, listings, clrscode3e, fancyvrb}
\def\code#1{\Verb!#1!}
\begin{document}
%begin
\begin {large}

\textbf{01.12.2023 Контрольна робота, варіант  \No { 1 }}\par


{
%beginitem
1. По яких днях тижня відбувались практичні з предмету?

%enditem


\vspace{5mm}
%beginitem
2. Як зовуть (ім'я, по батькові) викладача, що вів практичні заняття?

%enditem


\vspace{5mm}
%beginitem
3. Виберіть всі правильні твердження, якщо такі тут є.\par
\vspace{2mm}
( 1 ) Послідовність, що має експоненційну твірну, звичайної твірної може не мати.%good

( 2 ) Паросполучення можна промоделювати сполукою без повторень з обмеженнями на розташування елементів.%bad

( 3 ) Якщо послідовність має твірну функцію, то вона має і експоненційну твірну функцію.%good

( 4 ) Якщо дошка складається з кількох не зв’язаних між собою частин, то її туровий поліном можна отримати як добуток поліномів її частин.%bad

( 5 ) Якщо дошка складається з кількох частин, що не мають спільних рядків та спільних стовпчиків, то її туровий поліном можна отримати як добуток поліномів її частин.%good

( 6 ) Паросполучення можна промоделювати розміщенням з повтореннями з обмеженнями на розташування елементів.%bad

( 7 ) Експоненційна твірна функція послідовності будується виключно з використанням формули включення та виключення.%bad

( 8 ) Кожна послідовність має або ні твірну та експоненційну твірну одночасно.%bad

( 9 ) Задачу знаходження кількості способів складання кодових слів довжиною $n$, що складаються з десяткових цифр, де парні цифри зустрічаються непарну кількість разів, а парні цифри — непарну кількість разів, можна вирішити за допомогою експоненційної твірної.%good

( 10 ) Для дошки, що вкладається в квадрат ну дошку зі стороною 13, степінь турового полінома не перевищує 13.%good
%enditem
}
%end
\end {large}

\vspace{4mm}
%end
\newpage
%begin
\begin {large}

\textbf{01.12.2023 Контрольна робота, варіант  \No { 2 }}\par


{
%beginitem
1. По яких днях тижня відбувались практичні з предмету?

%enditem


\vspace{5mm}
%beginitem
2. Як зовуть (ім'я, по батькові) викладача, що вів практичні заняття?

%enditem


\vspace{5mm}
%beginitem
3. Виберіть всі правильні твердження, якщо такі тут є.\par
\vspace{2mm}
( 1 ) Послідовність, що має твірну функцію, експоненційної твірної може не мати.%bad

( 2 ) Задачу знаходження кількості способів сформувати новорічні подарункові набори з $n$ цукерок, що складаються з цукерок 13 видів, де цукерок кожного вигляду не більше 20, можна розв’язати методом твірних.мм%good

( 3 ) Паросполучення можна промоделювати розміщенням з повтореннями з обмеженнями на розташування елементів.%bad

( 4 ) Деякі задачі про перестановки з забороненими позиціями можуть бути розв’я\-за\-ні за допомогою формули включення та виключення.%good

( 5 ) Якщо послідовність має експоненційну твірну функцію, то вона має звичайну твірну функцію.%bad

( 6 ) Кожна числова послідовність має не більше однієї твірної.%good

( 7 ) Задачу знаходження кількості способів складання кодових слів довжиною $n$, що складаються з десяткових цифр, де парні цифри зустрічаються непарну кількість разів, а парні цифри — непарну кількість разів, можна вирішити за допомогою твірної.%bad

( 8 ) Кожна числова послідовність має не більше однієї  експоненційної твірної.%good

( 9 ) Послідовність може мати дві різних експоненційних твірні.%bad

( 10 ) Послідовність, задана рекурентним співвідношенням $$a_{n+1}=a_{n}a_{0}+a_{n-1}a_{1}+\ldots+a_0a_{n}$$ є послідовністю чисел Каталана.%bad
%enditem
}
%end
\end {large}

\vspace{4mm}
%end
\newpage
%begin
\begin {large}

\textbf{01.12.2023 Контрольна робота, варіант  \No { 3 }}\par


{
%beginitem
1. По яких днях тижня відбувались практичні з предмету?

%enditem


\vspace{5mm}
%beginitem
2. Як зовуть (ім'я, по батькові) викладача, що вів практичні заняття?

%enditem


\vspace{5mm}
%beginitem
3. Виберіть всі правильні твердження, якщо такі тут є.\par
\vspace{2mm}
( 1 ) Послідовність, задана рекурентним співвідношенням $a_{n+2}=a_{n+1}+a_{n}$ є послідовністю чисел Фібоначчі.%bad

( 2 ) Задачу знаходження кількості способів сформувати новорічні подарункові набори з $n$ цукерок, що складаються з цукерок 13 видів, де цукерок кожного вигляду не більше 20, можна розв’язати методом твірних.мм%good

( 3 ) Задачу знаходження кількості способів складання кодових слів довжиною $n$, що складаються з десяткових цифр, де парні цифри зустрічаються непарну кількість разів, а парні цифри — непарну кількість разів, можна вирішити за допомогою експоненційної твірної.%good

( 4 ) Формула включення та виключення будується за твірною функцією послідовності.%bad

( 5 ) Якщо послідовність має експоненційну твірну функцію, то вона має звичайну твірну функцію.%bad

( 6 ) Послідовність, задана рекурентним співвідношенням $$a_{n+1}=a_{n}a_{0}+a_{n-1}a_{1}+\ldots+a_0a_{n}$$ може не бути послідовністю чисел Каталана.%good

( 7 ) Послідовність, задана рекурентним співвідношенням $$a_{n+1}=a_{n}a_{0}+a_{n-1}a_{1}+\ldots+a_0a_{n}$$ є послідовністю чисел Каталана.%bad

( 8 ) Якщо дошка складається з кількох не зв’язаних між собою частин, то її туровий поліном можна отримати як суму поліномів її частин.%bad

( 9 ) Послідовність, задана рекурентним співвідношенням $a_{n+2}=a_{n+1}+a_{n}$ може не бути послідовністю чисел Фібоначчі.%good

( 10 ) Послідовність, задана рекурентним співвідношенням $$a_{n+1}=a_{n}a_{0}+a_{n-1}a_{1}+\ldots+a_0a_{n}$$ може не бути послідовністю чисел Каталана.%good
%enditem
}
%end
\end {large}

\vspace{4mm}
%end
\newpage
%begin
\begin {large}

\textbf{01.12.2023 Контрольна робота, варіант  \No { 4 }}\par


{
%beginitem
1. По яких днях тижня відбувались практичні з предмету?

%enditem


\vspace{5mm}
%beginitem
2. Як зовуть (ім'я, по батькові) викладача, що вів практичні заняття?

%enditem


\vspace{5mm}
%beginitem
3. Виберіть всі правильні твердження, якщо такі тут є.\par
\vspace{2mm}
( 1 ) За означенням туровий поліном є експоненційною твірною функцією.%bad

( 2 ) Послідовність, задана рекурентним співвідношенням $a_{n+2}=a_{n+1}+a_{n}$ може не бути послідовністю чисел Фібоначчі.%good

( 3 ) Для прямокутної дошки зі сторонами 5 та 7 степінь її турового полінома дорівнює 12.%bad

( 4 ) Задачу знаходження кількості способів сформувати новорічні подарункові набори з $n$ цукерок, що складаються з цукерок 13 видів, де цукерок кожного вигляду не більше 20, можна розв’язати за допомогою експоненційних твірних.%bad

( 5 ) Якщо дошка складається з кількох не зв’язаних між собою частин, то для неї турового полінома не існує.%bad

( 6 ) Задачу знаходження кількості способів сформувати новорічні подарункові набори з $n$ цукерок, що складаються з цукерок 13 видів, де цукерок кожного вигляду не більше 20, можна розв’язати методом рекурентних співвідношень.%good

( 7 ) Для прямокутної дошки зі сторонами 5 та 7 степінь її турового полінома дорівнює 12.%bad

( 8 ) Задачу знаходження кількості способів сформувати новорічні подарункові набори з $n$ цукерок, що складаються з цукерок 13 видів, де цукерок кожного вигляду не більше 20, можна розв’язати методом рекурентних співвідношень.%good

( 9 ) Для дошки, що вкладається в квадрат ну дошку зі стороною 13, степінь турового полінома не перевищує 13.%good

( 10 ) Задачу знаходження кількості способів складання кодових слів довжиною $n$, що складаються з десяткових цифр, де парні цифри зустрічаються непарну кількість разів, а парні цифри — непарну кількість разів, можна вирішити методом твірних.%good
%enditem
}
%end
\end {large}

\vspace{4mm}
%end
\newpage
%begin
\begin {large}

\textbf{01.12.2023 Контрольна робота, варіант  \No { 5 }}\par


{
%beginitem
1. По яких днях тижня відбувались практичні з предмету?

%enditem


\vspace{5mm}
%beginitem
2. Як зовуть (ім'я, по батькові) викладача, що вів практичні заняття?

%enditem


\vspace{5mm}
%beginitem
3. Виберіть всі правильні твердження, якщо такі тут є.\par
\vspace{2mm}
( 1 ) Паросполучення можна промоделювати сполукою без повторень з обмеженнями на розташування елементів.%bad

( 2 ) Питання, чи існує послідовність, що має дві різні твірні, є відкритим.%bad

( 3 ) Турові поліноми --- це ані про тури, ані про поліноми.%bad

( 4 ) Турові поліноми з шахами майже нічого спільного не мають.%good

( 5 ) Експоненційна твірна функція послідовності будується виключно з використанням формули включення та виключення.%bad

( 6 ) Якщо дошка складається з кількох не зв’язаних між собою частин, то для неї турового полінома не існує.%bad

( 7 ) Коли кажуть про твірну функцію послідовності, то на увазі завжди мають експоненційну твірну.%bad

( 8 ) Паросполучення можна промоделювати сполукою з повтореннями з обмеженнями на розташування елементів.%bad

( 9 ) Послідовність, що має експоненційну твірну, звичайної твірної може не мати.%good

( 10 ) Паросполучення можна промоделювати розміщенням з обмеженнями на розташування елементів.%good
%enditem
}
%end
\end {large}

\vspace{4mm}
%end
\newpage
%begin
\begin {large}

\textbf{01.12.2023 Контрольна робота, варіант  \No { 6 }}\par


{
%beginitem
1. По яких днях тижня відбувались практичні з предмету?

%enditem


\vspace{5mm}
%beginitem
2. Як зовуть (ім'я, по батькові) викладача, що вів практичні заняття?

%enditem


\vspace{5mm}
%beginitem
3. Виберіть всі правильні твердження, якщо такі тут є.\par
\vspace{2mm}
( 1 ) Турові поліноми можна розглядати не тільки для квадратних дошок.%good

( 2 ) За означенням туровий поліном є твірною функцією.%good

( 3 ) Якщо послідовність має твірну функцію, то вона має і експоненційну твірну функцію.%good

( 4 ) Для дошки, що вкладається в квадрат ну дошку зі стороною 13, степінь турового полінома в точності дорівнює 13.%bad

( 5 ) Задачу знаходження кількості способів складання кодових слів довжиною $n$, що складаються з десяткових цифр, де парні цифри зустрічаються непарну кількість разів, а парні цифри — непарну кількість разів, можна вирішити за допомогою твірної.%bad

( 6 ) Турові поліноми розглядаються виключно для квадратних дошок.%bad

( 7 ) Послідовність, що має твірну функцію, експоненційної твірної може не мати.%bad

( 8 ) Турові поліноми розглядаються виключно для квадратних дошок.%bad

( 9 ) Задачу знаходження кількості способів складання кодових слів довжиною $n$, що складаються з десяткових цифр, де парні цифри зустрічаються непарну кількість разів, а парні цифри — непарну кількість разів, можна вирішити за допомогою твірної.%bad

( 10 ) Питання, чи існує послідовність, що має дві різні твірні, є відкритим.%bad
%enditem
}
%end
\end {large}

\vspace{4mm}
%end
\newpage
%begin
\begin {large}

\textbf{01.12.2023 Контрольна робота, варіант  \No { 7 }}\par


{
%beginitem
1. По яких днях тижня відбувались практичні з предмету?

%enditem


\vspace{5mm}
%beginitem
2. Як зовуть (ім'я, по батькові) викладача, що вів практичні заняття?

%enditem


\vspace{5mm}
%beginitem
3. Виберіть всі правильні твердження, якщо такі тут є.\par
\vspace{2mm}
( 1 ) Турові поліноми можна розглядати не тільки для квадратних дошок.%good

( 2 ) Якщо дошка складається з кількох частин, що не мають спільних рядків та спільних стовпчиків, то її туровий поліном можна отримати як добуток поліномів її частин.%good

( 3 ) Кожна послідовність має або ні твірну та експоненційну твірну одночасно.%bad

( 4 ) Послідовність, що має твірну функцію, експоненційної твірної може не мати.%bad

( 5 ) Кожна послідовність має або ні твірну та експоненційну твірну одночасно.%bad

( 6 ) Послідовність, задана рекурентним співвідношенням $a_{n+2}=a_{n+1}+a_{n}$ є послідовністю чисел Фібоначчі.%bad

( 7 ) Задачу знаходження кількості способів складання кодових слів довжиною $n$, що складаються з десяткових цифр, де парні цифри зустрічаються непарну кількість разів, а парні цифри — непарну кількість разів, можна вирішити методом твірних.%good

( 8 ) Задачу знаходження кількості способів складання кодових слів довжиною $n$, що складаються з десяткових цифр, де парні цифри зустрічаються непарну кількість разів, а парні цифри — непарну кількість разів, можна вирішити методом твірних.%good

( 9 ) Послідовність, задана рекурентним співвідношенням $a_{n+2}=a_{n+1}+a_{n}$ є послідовністю чисел Фібоначчі.%bad

( 10 ) Питання, чи існує послідовність, що має дві різні твірні, є відкритим.%bad
%enditem
}
%end
\end {large}

\vspace{4mm}
%end
\newpage
%begin
\begin {large}

\textbf{01.12.2023 Контрольна робота, варіант  \No { 8 }}\par


{
%beginitem
1. По яких днях тижня відбувались практичні з предмету?

%enditem


\vspace{5mm}
%beginitem
2. Як зовуть (ім'я, по батькові) викладача, що вів практичні заняття?

%enditem


\vspace{5mm}
%beginitem
3. Виберіть всі правильні твердження, якщо такі тут є.\par
\vspace{2mm}
( 1 ) Послідовність, задана рекурентним співвідношенням $$a_{n+1}=a_{n}a_{0}+a_{n-1}a_{1}+\ldots+a_0a_{n}$$ може не бути послідовністю чисел Каталана.%good

( 2 ) Задачу знаходження кількості способів сформувати новорічні подарункові набори з $n$ цукерок, що складаються з цукерок 13 видів, де цукерок кожного вигляду не більше 20, можна розв’язати за допомогою експоненційних твірних.%bad

( 3 ) Послідовність може мати дві різних експоненційних твірні.%bad

( 4 ) Паросполучення можна промоделювати сполукою без повторень з обмеженнями на розташування елементів.%bad

( 5 ) За означенням туровий поліном є твірною функцією.%good

( 6 ) Формула включення та виключення будується за твірною функцією послідовності.%bad

( 7 ) Послідовність, задана рекурентним співвідношенням $a_{n+2}=a_{n+1}+a_{n}$ може не бути послідовністю чисел Фібоначчі.%good

( 8 ) Задачу знаходження кількості способів сформувати новорічні подарункові набори з $n$ цукерок, що складаються з цукерок 13 видів, де цукерок кожного вигляду не більше 20, можна розв’язати методом твірних.мм%good

( 9 ) За означенням туровий поліном є експоненційною твірною функцією.%bad

( 10 ) Турові поліноми з шахами майже нічого спільного не мають.%good
%enditem
}
%end
\end {large}

\vspace{4mm}
%end
\newpage
%begin
\begin {large}

\textbf{01.12.2023 Контрольна робота, варіант  \No { 9 }}\par


{
%beginitem
1. По яких днях тижня відбувались практичні з предмету?

%enditem


\vspace{5mm}
%beginitem
2. Як зовуть (ім'я, по батькові) викладача, що вів практичні заняття?

%enditem


\vspace{5mm}
%beginitem
3. Виберіть всі правильні твердження, якщо такі тут є.\par
\vspace{2mm}
( 1 ) Турові поліноми з шахами майже нічого спільного не мають.%good

( 2 ) Кожна числова послідовність має не більше однієї  експоненційної твірної.%good

( 3 ) Турові поліноми --- це ані про тури, ані про поліноми.%bad

( 4 ) Паросполучення можна промоделювати сполукою з повтореннями з обмеженнями на розташування елементів.%bad

( 5 ) Послідовність може мати дві різних експоненційних твірні.%bad

( 6 ) Якщо послідовність має експоненційну твірну функцію, то вона має звичайну твірну функцію.%bad

( 7 ) Коли кажуть про твірну функцію послідовності, то на увазі завжди мають експоненційну твірну.%bad

( 8 ) Послідовність, що має експоненційну твірну, звичайної твірної може не мати.%good

( 9 ) Якщо дошка складається з кількох не зв’язаних між собою частин, то її туровий поліном можна отримати як суму поліномів її частин.%bad

( 10 ) Задачу знаходження кількості способів сформувати новорічні подарункові набори з $n$ цукерок, що складаються з цукерок 13 видів, де цукерок кожного вигляду не більше 20, можна розв’язати за допомогою експоненційних твірних.%bad
%enditem
}
%end
\end {large}

\vspace{4mm}
%end
\newpage
%begin
\begin {large}

\textbf{01.12.2023 Контрольна робота, варіант  \No { 10 }}\par


{
%beginitem
1. По яких днях тижня відбувались практичні з предмету?

%enditem


\vspace{5mm}
%beginitem
2. Як зовуть (ім'я, по батькові) викладача, що вів практичні заняття?

%enditem


\vspace{5mm}
%beginitem
3. Виберіть всі правильні твердження, якщо такі тут є.\par
\vspace{2mm}
( 1 ) За означенням туровий поліном є експоненційною твірною функцією.%bad

( 2 ) Експоненційна твірна функція послідовності будується виключно з використанням формули включення та виключення.%bad

( 3 ) Коли кажуть про твірну функцію послідовності, то на увазі завжди мають експоненційну твірну.%bad

( 4 ) Кожна числова послідовність має не більше однієї  експоненційної твірної.%good

( 5 ) Послідовність, задана рекурентним співвідношенням $$a_{n+1}=a_{n}a_{0}+a_{n-1}a_{1}+\ldots+a_0a_{n}$$ є послідовністю чисел Каталана.%bad

( 6 ) Якщо дошка складається з кількох не зв’язаних між собою частин, то її туровий поліном можна отримати як добуток поліномів її частин.%bad

( 7 ) Паросполучення можна промоделювати розміщенням з повтореннями з обмеженнями на розташування елементів.%bad

( 8 ) Для дошки, що вкладається в квадрат ну дошку зі стороною 13, степінь турового полінома в точності дорівнює 13.%bad

( 9 ) Задачу знаходження кількості способів складання кодових слів довжиною $n$, що складаються з десяткових цифр, де парні цифри зустрічаються непарну кількість разів, а парні цифри — непарну кількість разів, можна вирішити за допомогою експоненційної твірної.%good

( 10 ) Для дошки, що вкладається в квадрат ну дошку зі стороною 13, степінь турового полінома не перевищує 13.%good
%enditem
}
%end
\end {large}

\vspace{4mm}
%end
\newpage
%begin
\begin {large}

\textbf{01.12.2023 Контрольна робота, варіант  \No { 11 }}\par


{
%beginitem
1. По яких днях тижня відбувались практичні з предмету?

%enditem


\vspace{5mm}
%beginitem
2. Як зовуть (ім'я, по батькові) викладача, що вів практичні заняття?

%enditem


\vspace{5mm}
%beginitem
3. Виберіть всі правильні твердження, якщо такі тут є.\par
\vspace{2mm}
( 1 ) Якщо дошка складається з кількох не зв’язаних між собою частин, то її туровий поліном можна отримати як добуток поліномів її частин.%bad

( 2 ) За означенням туровий поліном є твірною функцією.%good

( 3 ) Турові поліноми можна розглядати не тільки для квадратних дошок.%good

( 4 ) Паросполучення можна промоделювати розміщенням з обмеженнями на розташування елементів.%good

( 5 ) Деякі задачі про перестановки з забороненими позиціями можуть бути розв’я\-за\-ні за допомогою формули включення та виключення.%good

( 6 ) Формула включення та виключення будується за твірною функцією послідовності.%bad

( 7 ) Задачу знаходження кількості способів сформувати новорічні подарункові набори з $n$ цукерок, що складаються з цукерок 13 видів, де цукерок кожного вигляду не більше 20, можна розв’язати методом рекурентних співвідношень.%good

( 8 ) Турові поліноми розглядаються виключно для квадратних дошок.%bad

( 9 ) Кожна числова послідовність має не більше однієї твірної.%good

( 10 ) Якщо послідовність має твірну функцію, то вона має і експоненційну твірну функцію.%good
%enditem
}
%end
\end {large}

\vspace{4mm}
%end
\newpage
%begin
\begin {large}

\textbf{01.12.2023 Контрольна робота, варіант  \No { 12 }}\par


{
%beginitem
1. По яких днях тижня відбувались практичні з предмету?

%enditem


\vspace{5mm}
%beginitem
2. Як зовуть (ім'я, по батькові) викладача, що вів практичні заняття?

%enditem


\vspace{5mm}
%beginitem
3. Виберіть всі правильні твердження, якщо такі тут є.\par
\vspace{2mm}
( 1 ) Для дошки, що вкладається в квадрат ну дошку зі стороною 13, степінь турового полінома в точності дорівнює 13.%bad

( 2 ) Якщо дошка складається з кількох не зв’язаних між собою частин, то для неї турового полінома не існує.%bad

( 3 ) Для прямокутної дошки зі сторонами 5 та 7 степінь її турового полінома дорівнює 12.%bad

( 4 ) Кожна числова послідовність має не більше однієї твірної.%good

( 5 ) Послідовність, задана рекурентним співвідношенням $a_{n+2}=a_{n+1}+a_{n}$ може не бути послідовністю чисел Фібоначчі.%good

( 6 ) Турові поліноми --- це ані про тури, ані про поліноми.%bad

( 7 ) Якщо дошка складається з кількох частин, що не мають спільних рядків та спільних стовпчиків, то її туровий поліном можна отримати як добуток поліномів її частин.%good

( 8 ) Якщо послідовність має експоненційну твірну функцію, то вона має звичайну твірну функцію.%bad

( 9 ) Послідовність може мати дві різних експоненційних твірні.%bad

( 10 ) Турові поліноми можна розглядати не тільки для квадратних дошок.%good
%enditem
}
%end
\end {large}

\vspace{4mm}
%end
\newpage
%begin
\begin {large}

\textbf{01.12.2023 Контрольна робота, варіант  \No { 13 }}\par


{
%beginitem
1. По яких днях тижня відбувались практичні з предмету?

%enditem


\vspace{5mm}
%beginitem
2. Як зовуть (ім'я, по батькові) викладача, що вів практичні заняття?

%enditem


\vspace{5mm}
%beginitem
3. Виберіть всі правильні твердження, якщо такі тут є.\par
\vspace{2mm}
( 1 ) Паросполучення можна промоделювати розміщенням з обмеженнями на розташування елементів.%good

( 2 ) Задачу знаходження кількості способів сформувати новорічні подарункові набори з $n$ цукерок, що складаються з цукерок 13 видів, де цукерок кожного вигляду не більше 20, можна розв’язати методом рекурентних співвідношень.%good

( 3 ) Турові поліноми --- це ані про тури, ані про поліноми.%bad

( 4 ) Турові поліноми з шахами майже нічого спільного не мають.%good

( 5 ) Задачу знаходження кількості способів складання кодових слів довжиною $n$, що складаються з десяткових цифр, де парні цифри зустрічаються непарну кількість разів, а парні цифри — непарну кількість разів, можна вирішити методом твірних.%good

( 6 ) Паросполучення можна промоделювати сполукою з повтореннями з обмеженнями на розташування елементів.%bad

( 7 ) Формула включення та виключення будується за твірною функцією послідовності.%bad

( 8 ) Послідовність, задана рекурентним співвідношенням $$a_{n+1}=a_{n}a_{0}+a_{n-1}a_{1}+\ldots+a_0a_{n}$$ може не бути послідовністю чисел Каталана.%good

( 9 ) Деякі задачі про перестановки з забороненими позиціями можуть бути розв’я\-за\-ні за допомогою формули включення та виключення.%good

( 10 ) Деякі задачі про перестановки з забороненими позиціями можуть бути розв’я\-за\-ні за допомогою формули включення та виключення.%good
%enditem
}
%end
\end {large}

\vspace{4mm}
%end
\newpage
%begin
\begin {large}

\textbf{01.12.2023 Контрольна робота, варіант  \No { 14 }}\par


{
%beginitem
1. По яких днях тижня відбувались практичні з предмету?

%enditem


\vspace{5mm}
%beginitem
2. Як зовуть (ім'я, по батькові) викладача, що вів практичні заняття?

%enditem


\vspace{5mm}
%beginitem
3. Виберіть всі правильні твердження, якщо такі тут є.\par
\vspace{2mm}
( 1 ) За означенням туровий поліном є твірною функцією.%good

( 2 ) Якщо дошка складається з кількох не зв’язаних між собою частин, то її туровий поліном можна отримати як суму поліномів її частин.%bad

( 3 ) Паросполучення можна промоделювати сполукою з повтореннями з обмеженнями на розташування елементів.%bad

( 4 ) Послідовність може мати дві різних експоненційних твірні.%bad

( 5 ) Задачу знаходження кількості способів складання кодових слів довжиною $n$, що складаються з десяткових цифр, де парні цифри зустрічаються непарну кількість разів, а парні цифри — непарну кількість разів, можна вирішити за допомогою експоненційної твірної.%good

( 6 ) Для дошки, що вкладається в квадрат ну дошку зі стороною 13, степінь турового полінома в точності дорівнює 13.%bad

( 7 ) Формула включення та виключення будується за твірною функцією послідовності.%bad

( 8 ) Питання, чи існує послідовність, що має дві різні твірні, є відкритим.%bad

( 9 ) Якщо дошка складається з кількох частин, що не мають спільних рядків та спільних стовпчиків, то її туровий поліном можна отримати як добуток поліномів її частин.%good

( 10 ) Паросполучення можна промоделювати розміщенням з обмеженнями на розташування елементів.%good
%enditem
}
%end
\end {large}

\vspace{4mm}
%end
\newpage
%begin
\begin {large}

\textbf{01.12.2023 Контрольна робота, варіант  \No { 15 }}\par


{
%beginitem
1. По яких днях тижня відбувались практичні з предмету?

%enditem


\vspace{5mm}
%beginitem
2. Як зовуть (ім'я, по батькові) викладача, що вів практичні заняття?

%enditem


\vspace{5mm}
%beginitem
3. Виберіть всі правильні твердження, якщо такі тут є.\par
\vspace{2mm}
( 1 ) Для дошки, що вкладається в квадрат ну дошку зі стороною 13, степінь турового полінома не перевищує 13.%good

( 2 ) Задачу знаходження кількості способів складання кодових слів довжиною $n$, що складаються з десяткових цифр, де парні цифри зустрічаються непарну кількість разів, а парні цифри — непарну кількість разів, можна вирішити за допомогою твірної.%bad

( 3 ) Послідовність, задана рекурентним співвідношенням $$a_{n+1}=a_{n}a_{0}+a_{n-1}a_{1}+\ldots+a_0a_{n}$$ може не бути послідовністю чисел Каталана.%good

( 4 ) Послідовність, задана рекурентним співвідношенням $a_{n+2}=a_{n+1}+a_{n}$ є послідовністю чисел Фібоначчі.%bad

( 5 ) Послідовність, що має експоненційну твірну, звичайної твірної може не мати.%good

( 6 ) Кожна числова послідовність має не більше однієї  експоненційної твірної.%good

( 7 ) Послідовність, що має твірну функцію, експоненційної твірної може не мати.%bad

( 8 ) Кожна послідовність має або ні твірну та експоненційну твірну одночасно.%bad

( 9 ) Якщо послідовність має твірну функцію, то вона має і експоненційну твірну функцію.%good

( 10 ) Задачу знаходження кількості способів сформувати новорічні подарункові набори з $n$ цукерок, що складаються з цукерок 13 видів, де цукерок кожного вигляду не більше 20, можна розв’язати за допомогою експоненційних твірних.%bad
%enditem
}
%end
\end {large}

\vspace{4mm}
%end
\newpage
%begin
\begin {large}

\textbf{01.12.2023 Контрольна робота, варіант  \No { 16 }}\par


{
%beginitem
1. По яких днях тижня відбувались практичні з предмету?

%enditem


\vspace{5mm}
%beginitem
2. Як зовуть (ім'я, по батькові) викладача, що вів практичні заняття?

%enditem


\vspace{5mm}
%beginitem
3. Виберіть всі правильні твердження, якщо такі тут є.\par
\vspace{2mm}
( 1 ) Задачу знаходження кількості способів складання кодових слів довжиною $n$, що складаються з десяткових цифр, де парні цифри зустрічаються непарну кількість разів, а парні цифри — непарну кількість разів, можна вирішити за допомогою експоненційної твірної.%good

( 2 ) Експоненційна твірна функція послідовності будується виключно з використанням формули включення та виключення.%bad

( 3 ) Турові поліноми розглядаються виключно для квадратних дошок.%bad

( 4 ) Якщо дошка складається з кількох не зв’язаних між собою частин, то для неї турового полінома не існує.%bad

( 5 ) Задачу знаходження кількості способів сформувати новорічні подарункові набори з $n$ цукерок, що складаються з цукерок 13 видів, де цукерок кожного вигляду не більше 20, можна розв’язати методом рекурентних співвідношень.%good

( 6 ) Якщо послідовність має твірну функцію, то вона має і експоненційну твірну функцію.%good

( 7 ) Паросполучення можна промоделювати розміщенням з повтореннями з обмеженнями на розташування елементів.%bad

( 8 ) Якщо дошка складається з кількох не зв’язаних між собою частин, то її туровий поліном можна отримати як суму поліномів її частин.%bad

( 9 ) Паросполучення можна промоделювати розміщенням з обмеженнями на розташування елементів.%good

( 10 ) Задачу знаходження кількості способів сформувати новорічні подарункові набори з $n$ цукерок, що складаються з цукерок 13 видів, де цукерок кожного вигляду не більше 20, можна розв’язати методом твірних.мм%good
%enditem
}
%end
\end {large}

\vspace{4mm}
%end
\newpage
%begin
\begin {large}

\textbf{01.12.2023 Контрольна робота, варіант  \No { 17 }}\par


{
%beginitem
1. По яких днях тижня відбувались практичні з предмету?

%enditem


\vspace{5mm}
%beginitem
2. Як зовуть (ім'я, по батькові) викладача, що вів практичні заняття?

%enditem


\vspace{5mm}
%beginitem
3. Виберіть всі правильні твердження, якщо такі тут є.\par
\vspace{2mm}
( 1 ) Послідовність, що має експоненційну твірну, звичайної твірної може не мати.%good

( 2 ) Якщо дошка складається з кількох не зв’язаних між собою частин, то для неї турового полінома не існує.%bad

( 3 ) Якщо дошка складається з кількох не зв’язаних між собою частин, то її туровий поліном можна отримати як добуток поліномів її частин.%bad

( 4 ) Паросполучення можна промоделювати сполукою без повторень з обмеженнями на розташування елементів.%bad

( 5 ) Для дошки, що вкладається в квадрат ну дошку зі стороною 13, степінь турового полінома не перевищує 13.%good

( 6 ) Коли кажуть про твірну функцію послідовності, то на увазі завжди мають експоненційну твірну.%bad

( 7 ) За означенням туровий поліном є експоненційною твірною функцією.%bad

( 8 ) Якщо послідовність має експоненційну твірну функцію, то вона має звичайну твірну функцію.%bad

( 9 ) Для прямокутної дошки зі сторонами 5 та 7 степінь її турового полінома дорівнює 12.%bad

( 10 ) Послідовність, задана рекурентним співвідношенням $$a_{n+1}=a_{n}a_{0}+a_{n-1}a_{1}+\ldots+a_0a_{n}$$ є послідовністю чисел Каталана.%bad
%enditem
}
%end
\end {large}

\vspace{4mm}
%end
\newpage
%begin
\begin {large}

\textbf{01.12.2023 Контрольна робота, варіант  \No { 18 }}\par


{
%beginitem
1. По яких днях тижня відбувались практичні з предмету?

%enditem


\vspace{5mm}
%beginitem
2. Як зовуть (ім'я, по батькові) викладача, що вів практичні заняття?

%enditem


\vspace{5mm}
%beginitem
3. Виберіть всі правильні твердження, якщо такі тут є.\par
\vspace{2mm}
( 1 ) Кожна числова послідовність має не більше однієї твірної.%good

( 2 ) Кожна числова послідовність має не більше однієї  експоненційної твірної.%good

( 3 ) Турові поліноми з шахами майже нічого спільного не мають.%good

( 4 ) Задачу знаходження кількості способів складання кодових слів довжиною $n$, що складаються з десяткових цифр, де парні цифри зустрічаються непарну кількість разів, а парні цифри — непарну кількість разів, можна вирішити за допомогою твірної.%bad

( 5 ) Послідовність, задана рекурентним співвідношенням $$a_{n+1}=a_{n}a_{0}+a_{n-1}a_{1}+\ldots+a_0a_{n}$$ може не бути послідовністю чисел Каталана.%good

( 6 ) Експоненційна твірна функція послідовності будується виключно з використанням формули включення та виключення.%bad

( 7 ) Задачу знаходження кількості способів складання кодових слів довжиною $n$, що складаються з десяткових цифр, де парні цифри зустрічаються непарну кількість разів, а парні цифри — непарну кількість разів, можна вирішити за допомогою твірної.%bad

( 8 ) Послідовність, задана рекурентним співвідношенням $$a_{n+1}=a_{n}a_{0}+a_{n-1}a_{1}+\ldots+a_0a_{n}$$ є послідовністю чисел Каталана.%bad

( 9 ) Задачу знаходження кількості способів сформувати новорічні подарункові набори з $n$ цукерок, що складаються з цукерок 13 видів, де цукерок кожного вигляду не більше 20, можна розв’язати методом твірних.мм%good

( 10 ) Задачу знаходження кількості способів сформувати новорічні подарункові набори з $n$ цукерок, що складаються з цукерок 13 видів, де цукерок кожного вигляду не більше 20, можна розв’язати методом рекурентних співвідношень.%good
%enditem
}
%end
\end {large}

\vspace{4mm}
%end
\newpage
%begin
\begin {large}

\textbf{01.12.2023 Контрольна робота, варіант  \No { 19 }}\par


{
%beginitem
1. По яких днях тижня відбувались практичні з предмету?

%enditem


\vspace{5mm}
%beginitem
2. Як зовуть (ім'я, по батькові) викладача, що вів практичні заняття?

%enditem


\vspace{5mm}
%beginitem
3. Виберіть всі правильні твердження, якщо такі тут є.\par
\vspace{2mm}
( 1 ) Якщо дошка складається з кількох не зв’язаних між собою частин, то для неї турового полінома не існує.%bad

( 2 ) Паросполучення можна промоделювати розміщенням з обмеженнями на розташування елементів.%good

( 3 ) Для прямокутної дошки зі сторонами 5 та 7 степінь її турового полінома дорівнює 12.%bad

( 4 ) Для дошки, що вкладається в квадрат ну дошку зі стороною 13, степінь турового полінома в точності дорівнює 13.%bad

( 5 ) Послідовність, задана рекурентним співвідношенням $$a_{n+1}=a_{n}a_{0}+a_{n-1}a_{1}+\ldots+a_0a_{n}$$ є послідовністю чисел Каталана.%bad

( 6 ) Послідовність, задана рекурентним співвідношенням $a_{n+2}=a_{n+1}+a_{n}$ є послідовністю чисел Фібоначчі.%bad

( 7 ) Послідовність може мати дві різних експоненційних твірні.%bad

( 8 ) Для прямокутної дошки зі сторонами 5 та 7 степінь її турового полінома дорівнює 12.%bad

( 9 ) За означенням туровий поліном є твірною функцією.%good

( 10 ) Турові поліноми з шахами майже нічого спільного не мають.%good
%enditem
}
%end
\end {large}

\vspace{4mm}
%end
\newpage
%begin
\begin {large}

\textbf{01.12.2023 Контрольна робота, варіант  \No { 20 }}\par


{
%beginitem
1. По яких днях тижня відбувались практичні з предмету?

%enditem


\vspace{5mm}
%beginitem
2. Як зовуть (ім'я, по батькові) викладача, що вів практичні заняття?

%enditem


\vspace{5mm}
%beginitem
3. Виберіть всі правильні твердження, якщо такі тут є.\par
\vspace{2mm}
( 1 ) Якщо дошка складається з кількох частин, що не мають спільних рядків та спільних стовпчиків, то її туровий поліном можна отримати як добуток поліномів її частин.%good

( 2 ) Питання, чи існує послідовність, що має дві різні твірні, є відкритим.%bad

( 3 ) Якщо послідовність має експоненційну твірну функцію, то вона має звичайну твірну функцію.%bad

( 4 ) Турові поліноми --- це ані про тури, ані про поліноми.%bad

( 5 ) Коли кажуть про твірну функцію послідовності, то на увазі завжди мають експоненційну твірну.%bad

( 6 ) Для дошки, що вкладається в квадрат ну дошку зі стороною 13, степінь турового полінома в точності дорівнює 13.%bad

( 7 ) Якщо дошка складається з кількох не зв’язаних між собою частин, то її туровий поліном можна отримати як добуток поліномів її частин.%bad

( 8 ) Задачу знаходження кількості способів сформувати новорічні подарункові набори з $n$ цукерок, що складаються з цукерок 13 видів, де цукерок кожного вигляду не більше 20, можна розв’язати методом твірних.мм%good

( 9 ) Задачу знаходження кількості способів сформувати новорічні подарункові набори з $n$ цукерок, що складаються з цукерок 13 видів, де цукерок кожного вигляду не більше 20, можна розв’язати за допомогою експоненційних твірних.%bad

( 10 ) Формула включення та виключення будується за твірною функцією послідовності.%bad
%enditem
}
%end
\end {large}

\vspace{4mm}
%end
\newpage
%begin
\begin {large}

\textbf{01.12.2023 Контрольна робота, варіант  \No { 21 }}\par


{
%beginitem
1. По яких днях тижня відбувались практичні з предмету?

%enditem


\vspace{5mm}
%beginitem
2. Як зовуть (ім'я, по батькові) викладача, що вів практичні заняття?

%enditem


\vspace{5mm}
%beginitem
3. Виберіть всі правильні твердження, якщо такі тут є.\par
\vspace{2mm}
( 1 ) Коли кажуть про твірну функцію послідовності, то на увазі завжди мають експоненційну твірну.%bad

( 2 ) Турові поліноми розглядаються виключно для квадратних дошок.%bad

( 3 ) Паросполучення можна промоделювати сполукою без повторень з обмеженнями на розташування елементів.%bad

( 4 ) Послідовність, задана рекурентним співвідношенням $a_{n+2}=a_{n+1}+a_{n}$ може не бути послідовністю чисел Фібоначчі.%good

( 5 ) Послідовність, що має твірну функцію, експоненційної твірної може не мати.%bad

( 6 ) Питання, чи існує послідовність, що має дві різні твірні, є відкритим.%bad

( 7 ) Кожна послідовність має або ні твірну та експоненційну твірну одночасно.%bad

( 8 ) Для дошки, що вкладається в квадрат ну дошку зі стороною 13, степінь турового полінома не перевищує 13.%good

( 9 ) Кожна числова послідовність має не більше однієї  експоненційної твірної.%good

( 10 ) Задачу знаходження кількості способів складання кодових слів довжиною $n$, що складаються з десяткових цифр, де парні цифри зустрічаються непарну кількість разів, а парні цифри — непарну кількість разів, можна вирішити методом твірних.%good
%enditem
}
%end
\end {large}

\vspace{4mm}
%end
\newpage
%begin
\begin {large}

\textbf{01.12.2023 Контрольна робота, варіант  \No { 22 }}\par


{
%beginitem
1. По яких днях тижня відбувались практичні з предмету?

%enditem


\vspace{5mm}
%beginitem
2. Як зовуть (ім'я, по батькові) викладача, що вів практичні заняття?

%enditem


\vspace{5mm}
%beginitem
3. Виберіть всі правильні твердження, якщо такі тут є.\par
\vspace{2mm}
( 1 ) Деякі задачі про перестановки з забороненими позиціями можуть бути розв’я\-за\-ні за допомогою формули включення та виключення.%good

( 2 ) Якщо дошка складається з кількох частин, що не мають спільних рядків та спільних стовпчиків, то її туровий поліном можна отримати як добуток поліномів її частин.%good

( 3 ) За означенням туровий поліном є експоненційною твірною функцією.%bad

( 4 ) Послідовність, що має експоненційну твірну, звичайної твірної може не мати.%good

( 5 ) Турові поліноми --- це ані про тури, ані про поліноми.%bad

( 6 ) Задачу знаходження кількості способів сформувати новорічні подарункові набори з $n$ цукерок, що складаються з цукерок 13 видів, де цукерок кожного вигляду не більше 20, можна розв’язати за допомогою експоненційних твірних.%bad

( 7 ) Турові поліноми можна розглядати не тільки для квадратних дошок.%good

( 8 ) Послідовність, задана рекурентним співвідношенням $a_{n+2}=a_{n+1}+a_{n}$ може не бути послідовністю чисел Фібоначчі.%good

( 9 ) Турові поліноми розглядаються виключно для квадратних дошок.%bad

( 10 ) Кожна числова послідовність має не більше однієї твірної.%good
%enditem
}
%end
\end {large}

\vspace{4mm}
%end
\newpage
%begin
\begin {large}

\textbf{01.12.2023 Контрольна робота, варіант  \No { 23 }}\par


{
%beginitem
1. По яких днях тижня відбувались практичні з предмету?

%enditem


\vspace{5mm}
%beginitem
2. Як зовуть (ім'я, по батькові) викладача, що вів практичні заняття?

%enditem


\vspace{5mm}
%beginitem
3. Виберіть всі правильні твердження, якщо такі тут є.\par
\vspace{2mm}
( 1 ) Задачу знаходження кількості способів складання кодових слів довжиною $n$, що складаються з десяткових цифр, де парні цифри зустрічаються непарну кількість разів, а парні цифри — непарну кількість разів, можна вирішити методом твірних.%good

( 2 ) Паросполучення можна промоделювати сполукою з повтореннями з обмеженнями на розташування елементів.%bad

( 3 ) Задачу знаходження кількості способів складання кодових слів довжиною $n$, що складаються з десяткових цифр, де парні цифри зустрічаються непарну кількість разів, а парні цифри — непарну кількість разів, можна вирішити за допомогою експоненційної твірної.%good

( 4 ) За означенням туровий поліном є експоненційною твірною функцією.%bad

( 5 ) Якщо дошка складається з кількох не зв’язаних між собою частин, то її туровий поліном можна отримати як суму поліномів її частин.%bad

( 6 ) Послідовність, задана рекурентним співвідношенням $a_{n+2}=a_{n+1}+a_{n}$ є послідовністю чисел Фібоначчі.%bad

( 7 ) Якщо послідовність має твірну функцію, то вона має і експоненційну твірну функцію.%good

( 8 ) Паросполучення можна промоделювати розміщенням з повтореннями з обмеженнями на розташування елементів.%bad

( 9 ) Кожна послідовність має або ні твірну та експоненційну твірну одночасно.%bad

( 10 ) Питання, чи існує послідовність, що має дві різні твірні, є відкритим.%bad
%enditem
}
%end
\end {large}

\vspace{4mm}
%end
\newpage
%begin
\begin {large}

\textbf{01.12.2023 Контрольна робота, варіант  \No { 24 }}\par


{
%beginitem
1. По яких днях тижня відбувались практичні з предмету?

%enditem


\vspace{5mm}
%beginitem
2. Як зовуть (ім'я, по батькові) викладача, що вів практичні заняття?

%enditem


\vspace{5mm}
%beginitem
3. Виберіть всі правильні твердження, якщо такі тут є.\par
\vspace{2mm}
( 1 ) Паросполучення можна промоделювати розміщенням з повтореннями з обмеженнями на розташування елементів.%bad

( 2 ) Послідовність, що має твірну функцію, експоненційної твірної може не мати.%bad

( 3 ) Кожна числова послідовність має не більше однієї твірної.%good

( 4 ) Експоненційна твірна функція послідовності будується виключно з використанням формули включення та виключення.%bad

( 5 ) Кожна послідовність має або ні твірну та експоненційну твірну одночасно.%bad

( 6 ) Якщо дошка складається з кількох не зв’язаних між собою частин, то її туровий поліном можна отримати як добуток поліномів її частин.%bad

( 7 ) Якщо послідовність має твірну функцію, то вона має і експоненційну твірну функцію.%good

( 8 ) Якщо дошка складається з кількох не зв’язаних між собою частин, то для неї турового полінома не існує.%bad

( 9 ) Паросполучення можна промоделювати розміщенням з повтореннями з обмеженнями на розташування елементів.%bad

( 10 ) Для прямокутної дошки зі сторонами 5 та 7 степінь її турового полінома дорівнює 12.%bad
%enditem
}
%end
\end {large}

\vspace{4mm}
%end
\newpage
%begin
\begin {large}

\textbf{01.12.2023 Контрольна робота, варіант  \No { 25 }}\par


{
%beginitem
1. По яких днях тижня відбувались практичні з предмету?

%enditem


\vspace{5mm}
%beginitem
2. Як зовуть (ім'я, по батькові) викладача, що вів практичні заняття?

%enditem


\vspace{5mm}
%beginitem
3. Виберіть всі правильні твердження, якщо такі тут є.\par
\vspace{2mm}
( 1 ) Паросполучення можна промоделювати розміщенням з обмеженнями на розташування елементів.%good

( 2 ) За означенням туровий поліном є твірною функцією.%good

( 3 ) Якщо дошка складається з кількох не зв’язаних між собою частин, то її туровий поліном можна отримати як добуток поліномів її частин.%bad

( 4 ) Кожна числова послідовність має не більше однієї твірної.%good

( 5 ) Турові поліноми можна розглядати не тільки для квадратних дошок.%good

( 6 ) Деякі задачі про перестановки з забороненими позиціями можуть бути розв’я\-за\-ні за допомогою формули включення та виключення.%good

( 7 ) За означенням туровий поліном є твірною функцією.%good

( 8 ) Задачу знаходження кількості способів сформувати новорічні подарункові набори з $n$ цукерок, що складаються з цукерок 13 видів, де цукерок кожного вигляду не більше 20, можна розв’язати за допомогою експоненційних твірних.%bad

( 9 ) Якщо послідовність має експоненційну твірну функцію, то вона має звичайну твірну функцію.%bad

( 10 ) Паросполучення можна промоделювати сполукою з повтореннями з обмеженнями на розташування елементів.%bad
%enditem
}
%end
\end {large}

\vspace{4mm}
%end
\newpage
%begin
\begin {large}

\textbf{01.12.2023 Контрольна робота, варіант  \No { 26 }}\par


{
%beginitem
1. По яких днях тижня відбувались практичні з предмету?

%enditem


\vspace{5mm}
%beginitem
2. Як зовуть (ім'я, по батькові) викладача, що вів практичні заняття?

%enditem


\vspace{5mm}
%beginitem
3. Виберіть всі правильні твердження, якщо такі тут є.\par
\vspace{2mm}
( 1 ) Послідовність, задана рекурентним співвідношенням $$a_{n+1}=a_{n}a_{0}+a_{n-1}a_{1}+\ldots+a_0a_{n}$$ може не бути послідовністю чисел Каталана.%good

( 2 ) Паросполучення можна промоделювати сполукою без повторень з обмеженнями на розташування елементів.%bad

( 3 ) Турові поліноми з шахами майже нічого спільного не мають.%good

( 4 ) Коли кажуть про твірну функцію послідовності, то на увазі завжди мають експоненційну твірну.%bad

( 5 ) Послідовність, задана рекурентним співвідношенням $a_{n+2}=a_{n+1}+a_{n}$ є послідовністю чисел Фібоначчі.%bad

( 6 ) Задачу знаходження кількості способів складання кодових слів довжиною $n$, що складаються з десяткових цифр, де парні цифри зустрічаються непарну кількість разів, а парні цифри — непарну кількість разів, можна вирішити за допомогою експоненційної твірної.%good

( 7 ) Турові поліноми --- це ані про тури, ані про поліноми.%bad

( 8 ) Послідовність, задана рекурентним співвідношенням $a_{n+2}=a_{n+1}+a_{n}$ може не бути послідовністю чисел Фібоначчі.%good

( 9 ) Послідовність, що має твірну функцію, експоненційної твірної може не мати.%bad

( 10 ) Експоненційна твірна функція послідовності будується виключно з використанням формули включення та виключення.%bad
%enditem
}
%end
\end {large}

\vspace{4mm}
%end
\newpage
%begin
\begin {large}

\textbf{01.12.2023 Контрольна робота, варіант  \No { 27 }}\par


{
%beginitem
1. По яких днях тижня відбувались практичні з предмету?

%enditem


\vspace{5mm}
%beginitem
2. Як зовуть (ім'я, по батькові) викладача, що вів практичні заняття?

%enditem


\vspace{5mm}
%beginitem
3. Виберіть всі правильні твердження, якщо такі тут є.\par
\vspace{2mm}
( 1 ) Турові поліноми розглядаються виключно для квадратних дошок.%bad

( 2 ) Якщо дошка складається з кількох не зв’язаних між собою частин, то її туровий поліном можна отримати як суму поліномів її частин.%bad

( 3 ) Задачу знаходження кількості способів сформувати новорічні подарункові набори з $n$ цукерок, що складаються з цукерок 13 видів, де цукерок кожного вигляду не більше 20, можна розв’язати методом твірних.мм%good

( 4 ) Деякі задачі про перестановки з забороненими позиціями можуть бути розв’я\-за\-ні за допомогою формули включення та виключення.%good

( 5 ) Задачу знаходження кількості способів сформувати новорічні подарункові набори з $n$ цукерок, що складаються з цукерок 13 видів, де цукерок кожного вигляду не більше 20, можна розв’язати методом рекурентних співвідношень.%good

( 6 ) Турові поліноми можна розглядати не тільки для квадратних дошок.%good

( 7 ) Якщо дошка складається з кількох не зв’язаних між собою частин, то її туровий поліном можна отримати як суму поліномів її частин.%bad

( 8 ) Формула включення та виключення будується за твірною функцією послідовності.%bad

( 9 ) За означенням туровий поліном є експоненційною твірною функцією.%bad

( 10 ) Задачу знаходження кількості способів сформувати новорічні подарункові набори з $n$ цукерок, що складаються з цукерок 13 видів, де цукерок кожного вигляду не більше 20, можна розв’язати методом рекурентних співвідношень.%good
%enditem
}
%end
\end {large}

\vspace{4mm}
%end
\newpage
%begin
\begin {large}

\textbf{01.12.2023 Контрольна робота, варіант  \No { 28 }}\par


{
%beginitem
1. По яких днях тижня відбувались практичні з предмету?

%enditem


\vspace{5mm}
%beginitem
2. Як зовуть (ім'я, по батькові) викладача, що вів практичні заняття?

%enditem


\vspace{5mm}
%beginitem
3. Виберіть всі правильні твердження, якщо такі тут є.\par
\vspace{2mm}
( 1 ) Послідовність, задана рекурентним співвідношенням $$a_{n+1}=a_{n}a_{0}+a_{n-1}a_{1}+\ldots+a_0a_{n}$$ є послідовністю чисел Каталана.%bad

( 2 ) Для дошки, що вкладається в квадрат ну дошку зі стороною 13, степінь турового полінома не перевищує 13.%good

( 3 ) Послідовність, що має експоненційну твірну, звичайної твірної може не мати.%good

( 4 ) Послідовність може мати дві різних експоненційних твірні.%bad

( 5 ) Турові поліноми з шахами майже нічого спільного не мають.%good

( 6 ) Кожна числова послідовність має не більше однієї  експоненційної твірної.%good

( 7 ) Деякі задачі про перестановки з забороненими позиціями можуть бути розв’я\-за\-ні за допомогою формули включення та виключення.%good

( 8 ) За означенням туровий поліном є твірною функцією.%good

( 9 ) Паросполучення можна промоделювати сполукою з повтореннями з обмеженнями на розташування елементів.%bad

( 10 ) Для дошки, що вкладається в квадрат ну дошку зі стороною 13, степінь турового полінома в точності дорівнює 13.%bad
%enditem
}
%end
\end {large}

\vspace{4mm}
%end
\newpage
%begin
\begin {large}

\textbf{01.12.2023 Контрольна робота, варіант  \No { 29 }}\par


{
%beginitem
1. По яких днях тижня відбувались практичні з предмету?

%enditem


\vspace{5mm}
%beginitem
2. Як зовуть (ім'я, по батькові) викладача, що вів практичні заняття?

%enditem


\vspace{5mm}
%beginitem
3. Виберіть всі правильні твердження, якщо такі тут є.\par
\vspace{2mm}
( 1 ) Задачу знаходження кількості способів складання кодових слів довжиною $n$, що складаються з десяткових цифр, де парні цифри зустрічаються непарну кількість разів, а парні цифри — непарну кількість разів, можна вирішити методом твірних.%good

( 2 ) Для прямокутної дошки зі сторонами 5 та 7 степінь її турового полінома дорівнює 12.%bad

( 3 ) Якщо дошка складається з кількох частин, що не мають спільних рядків та спільних стовпчиків, то її туровий поліном можна отримати як добуток поліномів її частин.%good

( 4 ) Задачу знаходження кількості способів складання кодових слів довжиною $n$, що складаються з десяткових цифр, де парні цифри зустрічаються непарну кількість разів, а парні цифри — непарну кількість разів, можна вирішити за допомогою твірної.%bad

( 5 ) Якщо дошка складається з кількох не зв’язаних між собою частин, то для неї турового полінома не існує.%bad

( 6 ) Паросполучення можна промоделювати сполукою без повторень з обмеженнями на розташування елементів.%bad

( 7 ) Послідовність, задана рекурентним співвідношенням $a_{n+2}=a_{n+1}+a_{n}$ є послідовністю чисел Фібоначчі.%bad

( 8 ) Послідовність, задана рекурентним співвідношенням $a_{n+2}=a_{n+1}+a_{n}$ є послідовністю чисел Фібоначчі.%bad

( 9 ) Якщо послідовність має твірну функцію, то вона має і експоненційну твірну функцію.%good

( 10 ) Паросполучення можна промоделювати розміщенням з обмеженнями на розташування елементів.%good
%enditem
}
%end
\end {large}

\vspace{4mm}
%end
\newpage
%begin
\begin {large}

\textbf{01.12.2023 Контрольна робота, варіант  \No { 30 }}\par


{
%beginitem
1. По яких днях тижня відбувались практичні з предмету?

%enditem


\vspace{5mm}
%beginitem
2. Як зовуть (ім'я, по батькові) викладача, що вів практичні заняття?

%enditem


\vspace{5mm}
%beginitem
3. Виберіть всі правильні твердження, якщо такі тут є.\par
\vspace{2mm}
( 1 ) Послідовність може мати дві різних експоненційних твірні.%bad

( 2 ) Питання, чи існує послідовність, що має дві різні твірні, є відкритим.%bad

( 3 ) Паросполучення можна промоделювати сполукою з повтореннями з обмеженнями на розташування елементів.%bad

( 4 ) Задачу знаходження кількості способів сформувати новорічні подарункові набори з $n$ цукерок, що складаються з цукерок 13 видів, де цукерок кожного вигляду не більше 20, можна розв’язати методом рекурентних співвідношень.%good

( 5 ) Турові поліноми можна розглядати не тільки для квадратних дошок.%good

( 6 ) Послідовність, задана рекурентним співвідношенням $$a_{n+1}=a_{n}a_{0}+a_{n-1}a_{1}+\ldots+a_0a_{n}$$ може не бути послідовністю чисел Каталана.%good

( 7 ) Коли кажуть про твірну функцію послідовності, то на увазі завжди мають експоненційну твірну.%bad

( 8 ) Задачу знаходження кількості способів складання кодових слів довжиною $n$, що складаються з десяткових цифр, де парні цифри зустрічаються непарну кількість разів, а парні цифри — непарну кількість разів, можна вирішити за допомогою твірної.%bad

( 9 ) Паросполучення можна промоделювати розміщенням з обмеженнями на розташування елементів.%good

( 10 ) Послідовність, задана рекурентним співвідношенням $$a_{n+1}=a_{n}a_{0}+a_{n-1}a_{1}+\ldots+a_0a_{n}$$ є послідовністю чисел Каталана.%bad
%enditem
}
%end
\end {large}

\vspace{4mm}
%end
\newpage
%begin
\begin {large}

\textbf{01.12.2023 Контрольна робота, варіант  \No { 31 }}\par


{
%beginitem
1. По яких днях тижня відбувались практичні з предмету?

%enditem


\vspace{5mm}
%beginitem
2. Як зовуть (ім'я, по батькові) викладача, що вів практичні заняття?

%enditem


\vspace{5mm}
%beginitem
3. Виберіть всі правильні твердження, якщо такі тут є.\par
\vspace{2mm}
( 1 ) Задачу знаходження кількості способів сформувати новорічні подарункові набори з $n$ цукерок, що складаються з цукерок 13 видів, де цукерок кожного вигляду не більше 20, можна розв’язати за допомогою експоненційних твірних.%bad

( 2 ) Для прямокутної дошки зі сторонами 5 та 7 степінь її турового полінома дорівнює 12.%bad

( 3 ) Паросполучення можна промоделювати розміщенням з повтореннями з обмеженнями на розташування елементів.%bad

( 4 ) Задачу знаходження кількості способів сформувати новорічні подарункові набори з $n$ цукерок, що складаються з цукерок 13 видів, де цукерок кожного вигляду не більше 20, можна розв’язати методом твірних.мм%good

( 5 ) Кожна числова послідовність має не більше однієї  експоненційної твірної.%good

( 6 ) Якщо дошка складається з кількох не зв’язаних між собою частин, то її туровий поліном можна отримати як добуток поліномів її частин.%bad

( 7 ) Кожна послідовність має або ні твірну та експоненційну твірну одночасно.%bad

( 8 ) За означенням туровий поліном є експоненційною твірною функцією.%bad

( 9 ) Послідовність, задана рекурентним співвідношенням $a_{n+2}=a_{n+1}+a_{n}$ може не бути послідовністю чисел Фібоначчі.%good

( 10 ) Кожна числова послідовність має не більше однієї твірної.%good
%enditem
}
%end
\end {large}

\vspace{4mm}
%end
\newpage
%begin
\begin {large}

\textbf{01.12.2023 Контрольна робота, варіант  \No { 32 }}\par


{
%beginitem
1. По яких днях тижня відбувались практичні з предмету?

%enditem


\vspace{5mm}
%beginitem
2. Як зовуть (ім'я, по батькові) викладача, що вів практичні заняття?

%enditem


\vspace{5mm}
%beginitem
3. Виберіть всі правильні твердження, якщо такі тут є.\par
\vspace{2mm}
( 1 ) Послідовність, що має твірну функцію, експоненційної твірної може не мати.%bad

( 2 ) Турові поліноми з шахами майже нічого спільного не мають.%good

( 3 ) Якщо послідовність має експоненційну твірну функцію, то вона має звичайну твірну функцію.%bad

( 4 ) Задачу знаходження кількості способів сформувати новорічні подарункові набори з $n$ цукерок, що складаються з цукерок 13 видів, де цукерок кожного вигляду не більше 20, можна розв’язати методом твірних.мм%good

( 5 ) Якщо дошка складається з кількох частин, що не мають спільних рядків та спільних стовпчиків, то її туровий поліном можна отримати як добуток поліномів її частин.%good

( 6 ) Паросполучення можна промоделювати сполукою без повторень з обмеженнями на розташування елементів.%bad

( 7 ) Послідовність може мати дві різних експоненційних твірні.%bad

( 8 ) За означенням туровий поліном є твірною функцією.%good

( 9 ) Для дошки, що вкладається в квадрат ну дошку зі стороною 13, степінь турового полінома не перевищує 13.%good

( 10 ) Для дошки, що вкладається в квадрат ну дошку зі стороною 13, степінь турового полінома в точності дорівнює 13.%bad
%enditem
}
%end
\end {large}

\vspace{4mm}
%end
\newpage
%begin
\begin {large}

\textbf{01.12.2023 Контрольна робота, варіант  \No { 33 }}\par


{
%beginitem
1. По яких днях тижня відбувались практичні з предмету?

%enditem


\vspace{5mm}
%beginitem
2. Як зовуть (ім'я, по батькові) викладача, що вів практичні заняття?

%enditem


\vspace{5mm}
%beginitem
3. Виберіть всі правильні твердження, якщо такі тут є.\par
\vspace{2mm}
( 1 ) Якщо дошка складається з кількох не зв’язаних між собою частин, то її туровий поліном можна отримати як суму поліномів її частин.%bad

( 2 ) Задачу знаходження кількості способів складання кодових слів довжиною $n$, що складаються з десяткових цифр, де парні цифри зустрічаються непарну кількість разів, а парні цифри — непарну кількість разів, можна вирішити методом твірних.%good

( 3 ) Деякі задачі про перестановки з забороненими позиціями можуть бути розв’я\-за\-ні за допомогою формули включення та виключення.%good

( 4 ) Для дошки, що вкладається в квадрат ну дошку зі стороною 13, степінь турового полінома в точності дорівнює 13.%bad

( 5 ) Кожна послідовність має або ні твірну та експоненційну твірну одночасно.%bad

( 6 ) Послідовність, задана рекурентним співвідношенням $$a_{n+1}=a_{n}a_{0}+a_{n-1}a_{1}+\ldots+a_0a_{n}$$ може не бути послідовністю чисел Каталана.%good

( 7 ) Кожна числова послідовність має не більше однієї твірної.%good

( 8 ) Задачу знаходження кількості способів складання кодових слів довжиною $n$, що складаються з десяткових цифр, де парні цифри зустрічаються непарну кількість разів, а парні цифри — непарну кількість разів, можна вирішити за допомогою експоненційної твірної.%good

( 9 ) Турові поліноми --- це ані про тури, ані про поліноми.%bad

( 10 ) За означенням туровий поліном є експоненційною твірною функцією.%bad
%enditem
}
%end
\end {large}

\vspace{4mm}
%end
\newpage
%begin
\begin {large}

\textbf{01.12.2023 Контрольна робота, варіант  \No { 34 }}\par


{
%beginitem
1. По яких днях тижня відбувались практичні з предмету?

%enditem


\vspace{5mm}
%beginitem
2. Як зовуть (ім'я, по батькові) викладача, що вів практичні заняття?

%enditem


\vspace{5mm}
%beginitem
3. Виберіть всі правильні твердження, якщо такі тут є.\par
\vspace{2mm}
( 1 ) Турові поліноми розглядаються виключно для квадратних дошок.%bad

( 2 ) Кожна числова послідовність має не більше однієї  експоненційної твірної.%good

( 3 ) Якщо дошка складається з кількох не зв’язаних між собою частин, то для неї турового полінома не існує.%bad

( 4 ) Якщо послідовність має експоненційну твірну функцію, то вона має звичайну твірну функцію.%bad

( 5 ) Задачу знаходження кількості способів сформувати новорічні подарункові набори з $n$ цукерок, що складаються з цукерок 13 видів, де цукерок кожного вигляду не більше 20, можна розв’язати за допомогою експоненційних твірних.%bad

( 6 ) Паросполучення можна промоделювати розміщенням з повтореннями з обмеженнями на розташування елементів.%bad

( 7 ) Експоненційна твірна функція послідовності будується виключно з використанням формули включення та виключення.%bad

( 8 ) Якщо дошка складається з кількох не зв’язаних між собою частин, то її туровий поліном можна отримати як суму поліномів її частин.%bad

( 9 ) Експоненційна твірна функція послідовності будується виключно з використанням формули включення та виключення.%bad

( 10 ) Формула включення та виключення будується за твірною функцією послідовності.%bad
%enditem
}
%end
\end {large}

\vspace{4mm}
%end
\newpage
%begin
\begin {large}

\textbf{01.12.2023 Контрольна робота, варіант  \No { 35 }}\par


{
%beginitem
1. По яких днях тижня відбувались практичні з предмету?

%enditem


\vspace{5mm}
%beginitem
2. Як зовуть (ім'я, по батькові) викладача, що вів практичні заняття?

%enditem


\vspace{5mm}
%beginitem
3. Виберіть всі правильні твердження, якщо такі тут є.\par
\vspace{2mm}
( 1 ) Формула включення та виключення будується за твірною функцією послідовності.%bad

( 2 ) Послідовність, що має твірну функцію, експоненційної твірної може не мати.%bad

( 3 ) Послідовність, що має експоненційну твірну, звичайної твірної може не мати.%good

( 4 ) Послідовність може мати дві різних експоненційних твірні.%bad

( 5 ) Якщо дошка складається з кількох частин, що не мають спільних рядків та спільних стовпчиків, то її туровий поліном можна отримати як добуток поліномів її частин.%good

( 6 ) Послідовність, що має експоненційну твірну, звичайної твірної може не мати.%good

( 7 ) Якщо дошка складається з кількох не зв’язаних між собою частин, то її туровий поліном можна отримати як добуток поліномів її частин.%bad

( 8 ) Кожна послідовність має або ні твірну та експоненційну твірну одночасно.%bad

( 9 ) Для прямокутної дошки зі сторонами 5 та 7 степінь її турового полінома дорівнює 12.%bad

( 10 ) За означенням туровий поліном є твірною функцією.%good
%enditem
}
%end
\end {large}

\vspace{4mm}
%end
\newpage
%begin
\begin {large}

\textbf{01.12.2023 Контрольна робота, варіант  \No { 36 }}\par


{
%beginitem
1. По яких днях тижня відбувались практичні з предмету?

%enditem


\vspace{5mm}
%beginitem
2. Як зовуть (ім'я, по батькові) викладача, що вів практичні заняття?

%enditem


\vspace{5mm}
%beginitem
3. Виберіть всі правильні твердження, якщо такі тут є.\par
\vspace{2mm}
( 1 ) Якщо дошка складається з кількох частин, що не мають спільних рядків та спільних стовпчиків, то її туровий поліном можна отримати як добуток поліномів її частин.%good

( 2 ) Турові поліноми --- це ані про тури, ані про поліноми.%bad

( 3 ) Турові поліноми --- це ані про тури, ані про поліноми.%bad

( 4 ) Задачу знаходження кількості способів складання кодових слів довжиною $n$, що складаються з десяткових цифр, де парні цифри зустрічаються непарну кількість разів, а парні цифри — непарну кількість разів, можна вирішити методом твірних.%good

( 5 ) Якщо послідовність має експоненційну твірну функцію, то вона має звичайну твірну функцію.%bad

( 6 ) Якщо дошка складається з кількох не зв’язаних між собою частин, то її туровий поліном можна отримати як добуток поліномів її частин.%bad

( 7 ) Послідовність, задана рекурентним співвідношенням $$a_{n+1}=a_{n}a_{0}+a_{n-1}a_{1}+\ldots+a_0a_{n}$$ є послідовністю чисел Каталана.%bad

( 8 ) Задачу знаходження кількості способів сформувати новорічні подарункові набори з $n$ цукерок, що складаються з цукерок 13 видів, де цукерок кожного вигляду не більше 20, можна розв’язати методом твірних.мм%good

( 9 ) Послідовність, задана рекурентним співвідношенням $$a_{n+1}=a_{n}a_{0}+a_{n-1}a_{1}+\ldots+a_0a_{n}$$ може не бути послідовністю чисел Каталана.%good

( 10 ) Турові поліноми розглядаються виключно для квадратних дошок.%bad
%enditem
}
%end
\end {large}

\vspace{4mm}
%end
\newpage
%begin
\begin {large}

\textbf{01.12.2023 Контрольна робота, варіант  \No { 37 }}\par


{
%beginitem
1. По яких днях тижня відбувались практичні з предмету?

%enditem


\vspace{5mm}
%beginitem
2. Як зовуть (ім'я, по батькові) викладача, що вів практичні заняття?

%enditem


\vspace{5mm}
%beginitem
3. Виберіть всі правильні твердження, якщо такі тут є.\par
\vspace{2mm}
( 1 ) Кожна числова послідовність має не більше однієї твірної.%good

( 2 ) Паросполучення можна промоделювати розміщенням з обмеженнями на розташування елементів.%good

( 3 ) Послідовність, задана рекурентним співвідношенням $a_{n+2}=a_{n+1}+a_{n}$ може не бути послідовністю чисел Фібоначчі.%good

( 4 ) Турові поліноми розглядаються виключно для квадратних дошок.%bad

( 5 ) Послідовність, задана рекурентним співвідношенням $$a_{n+1}=a_{n}a_{0}+a_{n-1}a_{1}+\ldots+a_0a_{n}$$ є послідовністю чисел Каталана.%bad

( 6 ) Задачу знаходження кількості способів складання кодових слів довжиною $n$, що складаються з десяткових цифр, де парні цифри зустрічаються непарну кількість разів, а парні цифри — непарну кількість разів, можна вирішити за допомогою експоненційної твірної.%good

( 7 ) Формула включення та виключення будується за твірною функцією послідовності.%bad

( 8 ) Задачу знаходження кількості способів складання кодових слів довжиною $n$, що складаються з десяткових цифр, де парні цифри зустрічаються непарну кількість разів, а парні цифри — непарну кількість разів, можна вирішити методом твірних.%good

( 9 ) Якщо дошка складається з кількох не зв’язаних між собою частин, то її туровий поліном можна отримати як суму поліномів її частин.%bad

( 10 ) Турові поліноми з шахами майже нічого спільного не мають.%good
%enditem
}
%end
\end {large}

\vspace{4mm}
%end
\newpage
%begin
\begin {large}

\textbf{01.12.2023 Контрольна робота, варіант  \No { 38 }}\par


{
%beginitem
1. По яких днях тижня відбувались практичні з предмету?

%enditem


\vspace{5mm}
%beginitem
2. Як зовуть (ім'я, по батькові) викладача, що вів практичні заняття?

%enditem


\vspace{5mm}
%beginitem
3. Виберіть всі правильні твердження, якщо такі тут є.\par
\vspace{2mm}
( 1 ) Послідовність, задана рекурентним співвідношенням $a_{n+2}=a_{n+1}+a_{n}$ є послідовністю чисел Фібоначчі.%bad

( 2 ) Кожна числова послідовність має не більше однієї  експоненційної твірної.%good

( 3 ) Паросполучення можна промоделювати сполукою без повторень з обмеженнями на розташування елементів.%bad

( 4 ) Для дошки, що вкладається в квадрат ну дошку зі стороною 13, степінь турового полінома не перевищує 13.%good

( 5 ) Задачу знаходження кількості способів сформувати новорічні подарункові набори з $n$ цукерок, що складаються з цукерок 13 видів, де цукерок кожного вигляду не більше 20, можна розв’язати методом рекурентних співвідношень.%good

( 6 ) Послідовність, що має твірну функцію, експоненційної твірної може не мати.%bad

( 7 ) Задачу знаходження кількості способів складання кодових слів довжиною $n$, що складаються з десяткових цифр, де парні цифри зустрічаються непарну кількість разів, а парні цифри — непарну кількість разів, можна вирішити за допомогою твірної.%bad

( 8 ) Деякі задачі про перестановки з забороненими позиціями можуть бути розв’я\-за\-ні за допомогою формули включення та виключення.%good

( 9 ) Якщо послідовність має твірну функцію, то вона має і експоненційну твірну функцію.%good

( 10 ) Задачу знаходження кількості способів складання кодових слів довжиною $n$, що складаються з десяткових цифр, де парні цифри зустрічаються непарну кількість разів, а парні цифри — непарну кількість разів, можна вирішити за допомогою твірної.%bad
%enditem
}
%end
\end {large}

\vspace{4mm}
%end
\newpage
%begin
\begin {large}

\textbf{01.12.2023 Контрольна робота, варіант  \No { 39 }}\par


{
%beginitem
1. По яких днях тижня відбувались практичні з предмету?

%enditem


\vspace{5mm}
%beginitem
2. Як зовуть (ім'я, по батькові) викладача, що вів практичні заняття?

%enditem


\vspace{5mm}
%beginitem
3. Виберіть всі правильні твердження, якщо такі тут є.\par
\vspace{2mm}
( 1 ) Паросполучення можна промоделювати сполукою з повтореннями з обмеженнями на розташування елементів.%bad

( 2 ) Питання, чи існує послідовність, що має дві різні твірні, є відкритим.%bad

( 3 ) Турові поліноми можна розглядати не тільки для квадратних дошок.%good

( 4 ) Паросполучення можна промоделювати сполукою з повтореннями з обмеженнями на розташування елементів.%bad

( 5 ) Для дошки, що вкладається в квадрат ну дошку зі стороною 13, степінь турового полінома в точності дорівнює 13.%bad

( 6 ) Послідовність, задана рекурентним співвідношенням $a_{n+2}=a_{n+1}+a_{n}$ може не бути послідовністю чисел Фібоначчі.%good

( 7 ) Турові поліноми можна розглядати не тільки для квадратних дошок.%good

( 8 ) За означенням туровий поліном є експоненційною твірною функцією.%bad

( 9 ) Коли кажуть про твірну функцію послідовності, то на увазі завжди мають експоненційну твірну.%bad

( 10 ) Коли кажуть про твірну функцію послідовності, то на увазі завжди мають експоненційну твірну.%bad
%enditem
}
%end
\end {large}

\vspace{4mm}
%end
\newpage
%begin
\begin {large}

\textbf{01.12.2023 Контрольна робота, варіант  \No { 40 }}\par


{
%beginitem
1. По яких днях тижня відбувались практичні з предмету?

%enditem


\vspace{5mm}
%beginitem
2. Як зовуть (ім'я, по батькові) викладача, що вів практичні заняття?

%enditem


\vspace{5mm}
%beginitem
3. Виберіть всі правильні твердження, якщо такі тут є.\par
\vspace{2mm}
( 1 ) Експоненційна твірна функція послідовності будується виключно з використанням формули включення та виключення.%bad

( 2 ) Для дошки, що вкладається в квадрат ну дошку зі стороною 13, степінь турового полінома не перевищує 13.%good

( 3 ) Задачу знаходження кількості способів сформувати новорічні подарункові набори з $n$ цукерок, що складаються з цукерок 13 видів, де цукерок кожного вигляду не більше 20, можна розв’язати за допомогою експоненційних твірних.%bad

( 4 ) Паросполучення можна промоделювати сполукою без повторень з обмеженнями на розташування елементів.%bad

( 5 ) Послідовність, що має експоненційну твірну, звичайної твірної може не мати.%good

( 6 ) Паросполучення можна промоделювати розміщенням з повтореннями з обмеженнями на розташування елементів.%bad

( 7 ) Задачу знаходження кількості способів складання кодових слів довжиною $n$, що складаються з десяткових цифр, де парні цифри зустрічаються непарну кількість разів, а парні цифри — непарну кількість разів, можна вирішити за допомогою експоненційної твірної.%good

( 8 ) Якщо дошка складається з кількох не зв’язаних між собою частин, то для неї турового полінома не існує.%bad

( 9 ) Питання, чи існує послідовність, що має дві різні твірні, є відкритим.%bad

( 10 ) Якщо послідовність має твірну функцію, то вона має і експоненційну твірну функцію.%good
%enditem
}
%end
\end {large}

\vspace{4mm}
%end
\newpage
%begin
\begin {large}

\textbf{01.12.2023 Контрольна робота, варіант  \No { 41 }}\par


{
%beginitem
1. По яких днях тижня відбувались практичні з предмету?

%enditem


\vspace{5mm}
%beginitem
2. Як зовуть (ім'я, по батькові) викладача, що вів практичні заняття?

%enditem


\vspace{5mm}
%beginitem
3. Виберіть всі правильні твердження, якщо такі тут є.\par
\vspace{2mm}
( 1 ) Послідовність, задана рекурентним співвідношенням $a_{n+2}=a_{n+1}+a_{n}$ може не бути послідовністю чисел Фібоначчі.%good

( 2 ) Паросполучення можна промоделювати розміщенням з повтореннями з обмеженнями на розташування елементів.%bad

( 3 ) Паросполучення можна промоделювати розміщенням з обмеженнями на розташування елементів.%good

( 4 ) Задачу знаходження кількості способів сформувати новорічні подарункові набори з $n$ цукерок, що складаються з цукерок 13 видів, де цукерок кожного вигляду не більше 20, можна розв’язати методом твірних.мм%good

( 5 ) Задачу знаходження кількості способів складання кодових слів довжиною $n$, що складаються з десяткових цифр, де парні цифри зустрічаються непарну кількість разів, а парні цифри — непарну кількість разів, можна вирішити методом твірних.%good

( 6 ) Послідовність, задана рекурентним співвідношенням $$a_{n+1}=a_{n}a_{0}+a_{n-1}a_{1}+\ldots+a_0a_{n}$$ може не бути послідовністю чисел Каталана.%good

( 7 ) Задачу знаходження кількості способів складання кодових слів довжиною $n$, що складаються з десяткових цифр, де парні цифри зустрічаються непарну кількість разів, а парні цифри — непарну кількість разів, можна вирішити за допомогою твірної.%bad

( 8 ) Задачу знаходження кількості способів складання кодових слів довжиною $n$, що складаються з десяткових цифр, де парні цифри зустрічаються непарну кількість разів, а парні цифри — непарну кількість разів, можна вирішити за допомогою експоненційної твірної.%good

( 9 ) Задачу знаходження кількості способів сформувати новорічні подарункові набори з $n$ цукерок, що складаються з цукерок 13 видів, де цукерок кожного вигляду не більше 20, можна розв’язати методом рекурентних співвідношень.%good

( 10 ) Кожна послідовність має або ні твірну та експоненційну твірну одночасно.%bad
%enditem
}
%end
\end {large}

\vspace{4mm}
%end
\newpage
%begin
\begin {large}

\textbf{01.12.2023 Контрольна робота, варіант  \No { 42 }}\par


{
%beginitem
1. По яких днях тижня відбувались практичні з предмету?

%enditem


\vspace{5mm}
%beginitem
2. Як зовуть (ім'я, по батькові) викладача, що вів практичні заняття?

%enditem


\vspace{5mm}
%beginitem
3. Виберіть всі правильні твердження, якщо такі тут є.\par
\vspace{2mm}
( 1 ) Турові поліноми розглядаються виключно для квадратних дошок.%bad

( 2 ) Турові поліноми --- це ані про тури, ані про поліноми.%bad

( 3 ) Паросполучення можна промоделювати сполукою з повтореннями з обмеженнями на розташування елементів.%bad

( 4 ) Турові поліноми --- це ані про тури, ані про поліноми.%bad

( 5 ) Питання, чи існує послідовність, що має дві різні твірні, є відкритим.%bad

( 6 ) Для прямокутної дошки зі сторонами 5 та 7 степінь її турового полінома дорівнює 12.%bad

( 7 ) Формула включення та виключення будується за твірною функцією послідовності.%bad

( 8 ) Якщо дошка складається з кількох не зв’язаних між собою частин, то її туровий поліном можна отримати як суму поліномів її частин.%bad

( 9 ) Експоненційна твірна функція послідовності будується виключно з використанням формули включення та виключення.%bad

( 10 ) Якщо дошка складається з кількох не зв’язаних між собою частин, то її туровий поліном можна отримати як добуток поліномів її частин.%bad
%enditem
}
%end
\end {large}

\vspace{4mm}
%end
\newpage
%begin
\begin {large}

\textbf{01.12.2023 Контрольна робота, варіант  \No { 43 }}\par


{
%beginitem
1. По яких днях тижня відбувались практичні з предмету?

%enditem


\vspace{5mm}
%beginitem
2. Як зовуть (ім'я, по батькові) викладача, що вів практичні заняття?

%enditem


\vspace{5mm}
%beginitem
3. Виберіть всі правильні твердження, якщо такі тут є.\par
\vspace{2mm}
( 1 ) Якщо дошка складається з кількох не зв’язаних між собою частин, то для неї турового полінома не існує.%bad

( 2 ) Якщо послідовність має експоненційну твірну функцію, то вона має звичайну твірну функцію.%bad

( 3 ) Деякі задачі про перестановки з забороненими позиціями можуть бути розв’я\-за\-ні за допомогою формули включення та виключення.%good

( 4 ) Кожна числова послідовність має не більше однієї твірної.%good

( 5 ) Послідовність, задана рекурентним співвідношенням $a_{n+2}=a_{n+1}+a_{n}$ є послідовністю чисел Фібоначчі.%bad

( 6 ) Послідовність, що має твірну функцію, експоненційної твірної може не мати.%bad

( 7 ) Якщо дошка складається з кількох не зв’язаних між собою частин, то її туровий поліном можна отримати як добуток поліномів її частин.%bad

( 8 ) Турові поліноми можна розглядати не тільки для квадратних дошок.%good

( 9 ) Паросполучення можна промоделювати сполукою без повторень з обмеженнями на розташування елементів.%bad

( 10 ) За означенням туровий поліном є експоненційною твірною функцією.%bad
%enditem
}
%end
\end {large}

\vspace{4mm}
%end
\newpage
%begin
\begin {large}

\textbf{01.12.2023 Контрольна робота, варіант  \No { 44 }}\par


{
%beginitem
1. По яких днях тижня відбувались практичні з предмету?

%enditem


\vspace{5mm}
%beginitem
2. Як зовуть (ім'я, по батькові) викладача, що вів практичні заняття?

%enditem


\vspace{5mm}
%beginitem
3. Виберіть всі правильні твердження, якщо такі тут є.\par
\vspace{2mm}
( 1 ) Коли кажуть про твірну функцію послідовності, то на увазі завжди мають експоненційну твірну.%bad

( 2 ) Задачу знаходження кількості способів сформувати новорічні подарункові набори з $n$ цукерок, що складаються з цукерок 13 видів, де цукерок кожного вигляду не більше 20, можна розв’язати за допомогою експоненційних твірних.%bad

( 3 ) Послідовність, що має твірну функцію, експоненційної твірної може не мати.%bad

( 4 ) Турові поліноми з шахами майже нічого спільного не мають.%good

( 5 ) Задачу знаходження кількості способів складання кодових слів довжиною $n$, що складаються з десяткових цифр, де парні цифри зустрічаються непарну кількість разів, а парні цифри — непарну кількість разів, можна вирішити за допомогою твірної.%bad

( 6 ) Для дошки, що вкладається в квадрат ну дошку зі стороною 13, степінь турового полінома в точності дорівнює 13.%bad

( 7 ) Послідовність, задана рекурентним співвідношенням $$a_{n+1}=a_{n}a_{0}+a_{n-1}a_{1}+\ldots+a_0a_{n}$$ є послідовністю чисел Каталана.%bad

( 8 ) Турові поліноми можна розглядати не тільки для квадратних дошок.%good

( 9 ) За означенням туровий поліном є твірною функцією.%good

( 10 ) Кожна послідовність має або ні твірну та експоненційну твірну одночасно.%bad
%enditem
}
%end
\end {large}

\vspace{4mm}
%end
\newpage
%begin
\begin {large}

\textbf{01.12.2023 Контрольна робота, варіант  \No { 45 }}\par


{
%beginitem
1. По яких днях тижня відбувались практичні з предмету?

%enditem


\vspace{5mm}
%beginitem
2. Як зовуть (ім'я, по батькові) викладача, що вів практичні заняття?

%enditem


\vspace{5mm}
%beginitem
3. Виберіть всі правильні твердження, якщо такі тут є.\par
\vspace{2mm}
( 1 ) Якщо послідовність має твірну функцію, то вона має і експоненційну твірну функцію.%good

( 2 ) Для дошки, що вкладається в квадрат ну дошку зі стороною 13, степінь турового полінома не перевищує 13.%good

( 3 ) Турові поліноми з шахами майже нічого спільного не мають.%good

( 4 ) Якщо дошка складається з кількох частин, що не мають спільних рядків та спільних стовпчиків, то її туровий поліном можна отримати як добуток поліномів її частин.%good

( 5 ) Задачу знаходження кількості способів сформувати новорічні подарункові набори з $n$ цукерок, що складаються з цукерок 13 видів, де цукерок кожного вигляду не більше 20, можна розв’язати методом рекурентних співвідношень.%good

( 6 ) Турові поліноми розглядаються виключно для квадратних дошок.%bad

( 7 ) Якщо дошка складається з кількох частин, що не мають спільних рядків та спільних стовпчиків, то її туровий поліном можна отримати як добуток поліномів її частин.%good

( 8 ) Кожна числова послідовність має не більше однієї  експоненційної твірної.%good

( 9 ) Послідовність може мати дві різних експоненційних твірні.%bad

( 10 ) Послідовність, задана рекурентним співвідношенням $a_{n+2}=a_{n+1}+a_{n}$ є послідовністю чисел Фібоначчі.%bad
%enditem
}
%end
\end {large}

\vspace{4mm}
%end
\newpage
%begin
\begin {large}

\textbf{01.12.2023 Контрольна робота, варіант  \No { 46 }}\par


{
%beginitem
1. По яких днях тижня відбувались практичні з предмету?

%enditem


\vspace{5mm}
%beginitem
2. Як зовуть (ім'я, по батькові) викладача, що вів практичні заняття?

%enditem


\vspace{5mm}
%beginitem
3. Виберіть всі правильні твердження, якщо такі тут є.\par
\vspace{2mm}
( 1 ) Послідовність, що має експоненційну твірну, звичайної твірної може не мати.%good

( 2 ) Задачу знаходження кількості способів складання кодових слів довжиною $n$, що складаються з десяткових цифр, де парні цифри зустрічаються непарну кількість разів, а парні цифри — непарну кількість разів, можна вирішити за допомогою експоненційної твірної.%good

( 3 ) Експоненційна твірна функція послідовності будується виключно з використанням формули включення та виключення.%bad

( 4 ) Для дошки, що вкладається в квадрат ну дошку зі стороною 13, степінь турового полінома не перевищує 13.%good

( 5 ) Якщо дошка складається з кількох не зв’язаних між собою частин, то її туровий поліном можна отримати як суму поліномів її частин.%bad

( 6 ) Задачу знаходження кількості способів сформувати новорічні подарункові набори з $n$ цукерок, що складаються з цукерок 13 видів, де цукерок кожного вигляду не більше 20, можна розв’язати методом твірних.мм%good

( 7 ) Коли кажуть про твірну функцію послідовності, то на увазі завжди мають експоненційну твірну.%bad

( 8 ) Задачу знаходження кількості способів сформувати новорічні подарункові набори з $n$ цукерок, що складаються з цукерок 13 видів, де цукерок кожного вигляду не більше 20, можна розв’язати методом твірних.мм%good

( 9 ) Якщо дошка складається з кількох частин, що не мають спільних рядків та спільних стовпчиків, то її туровий поліном можна отримати як добуток поліномів її частин.%good

( 10 ) Формула включення та виключення будується за твірною функцією послідовності.%bad
%enditem
}
%end
\end {large}

\vspace{4mm}
%end
\newpage
%begin
\begin {large}

\textbf{01.12.2023 Контрольна робота, варіант  \No { 47 }}\par


{
%beginitem
1. По яких днях тижня відбувались практичні з предмету?

%enditem


\vspace{5mm}
%beginitem
2. Як зовуть (ім'я, по батькові) викладача, що вів практичні заняття?

%enditem


\vspace{5mm}
%beginitem
3. Виберіть всі правильні твердження, якщо такі тут є.\par
\vspace{2mm}
( 1 ) Якщо послідовність має твірну функцію, то вона має і експоненційну твірну функцію.%good

( 2 ) Для дошки, що вкладається в квадрат ну дошку зі стороною 13, степінь турового полінома в точності дорівнює 13.%bad

( 3 ) Паросполучення можна промоделювати сполукою з повтореннями з обмеженнями на розташування елементів.%bad

( 4 ) Питання, чи існує послідовність, що має дві різні твірні, є відкритим.%bad

( 5 ) Кожна числова послідовність має не більше однієї  експоненційної твірної.%good

( 6 ) Кожна послідовність має або ні твірну та експоненційну твірну одночасно.%bad

( 7 ) Паросполучення можна промоделювати розміщенням з повтореннями з обмеженнями на розташування елементів.%bad

( 8 ) Для дошки, що вкладається в квадрат ну дошку зі стороною 13, степінь турового полінома в точності дорівнює 13.%bad

( 9 ) За означенням туровий поліном є твірною функцією.%good

( 10 ) Формула включення та виключення будується за твірною функцією послідовності.%bad
%enditem
}
%end
\end {large}

\vspace{4mm}
%end
\newpage
%begin
\begin {large}

\textbf{01.12.2023 Контрольна робота, варіант  \No { 48 }}\par


{
%beginitem
1. По яких днях тижня відбувались практичні з предмету?

%enditem


\vspace{5mm}
%beginitem
2. Як зовуть (ім'я, по батькові) викладача, що вів практичні заняття?

%enditem


\vspace{5mm}
%beginitem
3. Виберіть всі правильні твердження, якщо такі тут є.\par
\vspace{2mm}
( 1 ) Послідовність може мати дві різних експоненційних твірні.%bad

( 2 ) Коли кажуть про твірну функцію послідовності, то на увазі завжди мають експоненційну твірну.%bad

( 3 ) Якщо послідовність має експоненційну твірну функцію, то вона має звичайну твірну функцію.%bad

( 4 ) Якщо послідовність має експоненційну твірну функцію, то вона має звичайну твірну функцію.%bad

( 5 ) За означенням туровий поліном є експоненційною твірною функцією.%bad

( 6 ) Послідовність, задана рекурентним співвідношенням $$a_{n+1}=a_{n}a_{0}+a_{n-1}a_{1}+\ldots+a_0a_{n}$$ може не бути послідовністю чисел Каталана.%good

( 7 ) Паросполучення можна промоделювати розміщенням з повтореннями з обмеженнями на розташування елементів.%bad

( 8 ) Послідовність, що має експоненційну твірну, звичайної твірної може не мати.%good

( 9 ) Турові поліноми розглядаються виключно для квадратних дошок.%bad

( 10 ) Кожна числова послідовність має не більше однієї  експоненційної твірної.%good
%enditem
}
%end
\end {large}

\vspace{4mm}
%end
\newpage
%begin
\begin {large}

\textbf{01.12.2023 Контрольна робота, варіант  \No { 49 }}\par


{
%beginitem
1. По яких днях тижня відбувались практичні з предмету?

%enditem


\vspace{5mm}
%beginitem
2. Як зовуть (ім'я, по батькові) викладача, що вів практичні заняття?

%enditem


\vspace{5mm}
%beginitem
3. Виберіть всі правильні твердження, якщо такі тут є.\par
\vspace{2mm}
( 1 ) Питання, чи існує послідовність, що має дві різні твірні, є відкритим.%bad

( 2 ) Для прямокутної дошки зі сторонами 5 та 7 степінь її турового полінома дорівнює 12.%bad

( 3 ) Паросполучення можна промоделювати розміщенням з обмеженнями на розташування елементів.%good

( 4 ) Для дошки, що вкладається в квадрат ну дошку зі стороною 13, степінь турового полінома не перевищує 13.%good

( 5 ) Якщо дошка складається з кількох не зв’язаних між собою частин, то її туровий поліном можна отримати як суму поліномів її частин.%bad

( 6 ) Послідовність, задана рекурентним співвідношенням $$a_{n+1}=a_{n}a_{0}+a_{n-1}a_{1}+\ldots+a_0a_{n}$$ може не бути послідовністю чисел Каталана.%good

( 7 ) Задачу знаходження кількості способів складання кодових слів довжиною $n$, що складаються з десяткових цифр, де парні цифри зустрічаються непарну кількість разів, а парні цифри — непарну кількість разів, можна вирішити методом твірних.%good

( 8 ) Задачу знаходження кількості способів складання кодових слів довжиною $n$, що складаються з десяткових цифр, де парні цифри зустрічаються непарну кількість разів, а парні цифри — непарну кількість разів, можна вирішити за допомогою експоненційної твірної.%good

( 9 ) Послідовність, задана рекурентним співвідношенням $$a_{n+1}=a_{n}a_{0}+a_{n-1}a_{1}+\ldots+a_0a_{n}$$ є послідовністю чисел Каталана.%bad

( 10 ) Задачу знаходження кількості способів складання кодових слів довжиною $n$, що складаються з десяткових цифр, де парні цифри зустрічаються непарну кількість разів, а парні цифри — непарну кількість разів, можна вирішити методом твірних.%good
%enditem
}
%end
\end {large}

\vspace{4mm}
%end
\newpage
%begin
\begin {large}

\textbf{01.12.2023 Контрольна робота, варіант  \No { 50 }}\par


{
%beginitem
1. По яких днях тижня відбувались практичні з предмету?

%enditem


\vspace{5mm}
%beginitem
2. Як зовуть (ім'я, по батькові) викладача, що вів практичні заняття?

%enditem


\vspace{5mm}
%beginitem
3. Виберіть всі правильні твердження, якщо такі тут є.\par
\vspace{2mm}
( 1 ) Послідовність, задана рекурентним співвідношенням $a_{n+2}=a_{n+1}+a_{n}$ може не бути послідовністю чисел Фібоначчі.%good

( 2 ) За означенням туровий поліном є твірною функцією.%good

( 3 ) Для прямокутної дошки зі сторонами 5 та 7 степінь її турового полінома дорівнює 12.%bad

( 4 ) Кожна числова послідовність має не більше однієї твірної.%good

( 5 ) Турові поліноми --- це ані про тури, ані про поліноми.%bad

( 6 ) Експоненційна твірна функція послідовності будується виключно з використанням формули включення та виключення.%bad

( 7 ) Задачу знаходження кількості способів сформувати новорічні подарункові набори з $n$ цукерок, що складаються з цукерок 13 видів, де цукерок кожного вигляду не більше 20, можна розв’язати за допомогою експоненційних твірних.%bad

( 8 ) Послідовність, задана рекурентним співвідношенням $a_{n+2}=a_{n+1}+a_{n}$ є послідовністю чисел Фібоначчі.%bad

( 9 ) Задачу знаходження кількості способів складання кодових слів довжиною $n$, що складаються з десяткових цифр, де парні цифри зустрічаються непарну кількість разів, а парні цифри — непарну кількість разів, можна вирішити за допомогою твірної.%bad

( 10 ) Якщо дошка складається з кількох не зв’язаних між собою частин, то для неї турового полінома не існує.%bad
%enditem
}
%end
\end {large}

\vspace{4mm}
%end
\newpage
\end{document}

